\documentclass[11pt,a4paper]{article}
\usepackage[T1]{fontenc}
\usepackage{lmodern}
\usepackage[margin=27mm]{geometry}
\usepackage{amsmath,amssymb,amsthm,microtype,booktabs}
\usepackage[hidelinks]{hyperref}
\newtheorem{theorem}{Theorem}
\newtheorem{lemma}[theorem]{Lemma}
\newtheorem{proposition}[theorem]{Proposition}
\newcommand{\Ostar}{O^{*}}
\newcommand{\E}{\mathbb E}
\newcommand{\Prb}{\mathbb P}
\emergencystretch=2em
\clubpenalty=10000
\widowpenalty=10000
\title{Replacing Random Signs by Neighborhood Counts\\
\large A deterministic early-enumeration theorem and a limitation}
\author{}
\date{8 October 2026}
\begin{document}
\maketitle

\begin{abstract}
The cancellation needed for early enumeration has an exact formulation
in terms of agreement and disagreement counts between signed
neighborhoods. For a fixed family of learned inequalities, these pair
counts determine the variance of their sum under a uniform completion;
independence between neighborhoods is unnecessary. We formulate a general
deterministic theorem using signed co-occurrence counts and counts of
residual-gap profiles. A bipartite special case illustrates the mechanism,
but does not improve on the simpler algorithm available for bipartite
formulas. We also show why pairwise balance alone is insufficient
for this particular summed-constraint generator: with zero initial gaps,
it can enumerate a constant fraction of all assignments despite exactly
orthogonal signed neighborhoods. These are sufficient conditions and a
limitation of one algorithm, not a characterization of hard Max 2-SAT.
\end{abstract}

\section{What independence means here}

Independence in the random-sign model is a property of the joint
distribution over formulas. Empirical counts inside one formula are
different objects. Even exactly balanced one- and two-coordinate counts
do not imply joint independence. For example, the three sign patterns
$(u,v,uv)$ over the four choices $(u,v)\in\{-1,+1\}^2$ have uniform
one- and two-coordinate frequencies but only four of eight possible
triples.

We should replace independence by the particular counting inequalities
used by the algorithm. This note separates two jobs: making learned
inequalities sufficiently restrictive, and preventing their bad values
from aligning too strongly. Pair counts handle the second job directly.
In the theorem below, literal-degree gaps supply the first job.

\section{Exact agreement counts for a fixed constraint family}

Consider fixed inequalities $L_i(y)\le k_i$ on $r$ free variables,
each counting one bad value on its support $N_i$. Write assignments as
$y_w\in\{-1,+1\}$ and set
\[
 b_i(w)=\begin{cases}
  +1,&\text{if value $+1$ is bad for $i$ at $w$},\\
  -1,&\text{if value $-1$ is bad for $i$ at $w$},\\
  0,&w\notin N_i.
 \end{cases}
\]
Thus $L_i(y)=|N_i|/2+\frac12\sum_w b_i(w)y_w$.
For two rows define the directly countable quantity
\begin{equation}
 C_{ij}=A_{ij}-O_{ij}=\sum_w b_i(w)b_j(w),             \label{eq:C}
\end{equation}
where $A_{ij}$ counts shared variables with agreeing bad values and
$O_{ij}$ counts shared variables with opposite bad values.
In particular, $C_{ii}=|N_i|$.

For a selected family $I$, define
\begin{align}
 s_I&=\sum_{i\in I}\left(\frac{|N_i|}{2}-k_i\right),\\
 V_I&=\sum_w\left(\sum_{i\in I}b_i(w)\right)^2
     =\sum_{i\in I}|N_i|+
                       2\sum_{\substack{i<j\\i,j\in I}}C_{ij}.
                                                        \label{eq:V}
\end{align}
The sum of the inequalities is
$Q_I(y)\le-s_I$, where
$Q_I(y)=\frac12\sum_w(\sum_{i\in I}b_i(w))y_w$.

\begin{lemma}[A deterministic counting bound]\label{lem:count}
If $s_I>0$ and $V_I>0$, the number of assignments satisfying the
summed inequality is at most
\begin{equation}
 2^r\exp\left(-\frac{2s_I^2}{V_I}\right).           \label{eq:tail}
\end{equation}
If $s_I>0$ and $V_I=0$, that number is zero.
\end{lemma}
\begin{proof}
Use a uniform assignment only to count completions of this fixed family.
Its coordinates $y_w$ are independent signs, regardless of any
dependencies between the rows $b_i$. Thus
\[
 \E e^{-\lambda Q_I}
 =\prod_w\cosh\left(\frac\lambda2\sum_{i\in I}b_i(w)\right)
 \le e^{\lambda^2V_I/8}.
\]
Markov's inequality with $\lambda=4s_I/V_I$ gives
\eqref{eq:tail}. If $V_I=0$, then $Q_I$ is identically zero.
\end{proof}

This is stronger than an analogy with independence: the exact variance
is $V_I/4$, and equation~\eqref{eq:V} expresses it using pair counts.
Independence of the neighborhoods is not used. For example, when every
$C_{ij}=0$ for $i\ne j$, then $V_I=\sum_{i\in I}|N_i|$ for every
selected subfamily, even if there are higher-order relations among rows.

For exact counts, let $m_w^+$ and $m_w^-$ count bad occurrences at $w$,
put $c_w=|m_w^+-m_w^-|$, and
$B=\sum_i k_i-\sum_w\min(m_w^+,m_w^-)$. Orient the cheaper value as zero.
The summed inequality is $\sum_wc_wz_w\le B$, and its exact completion
count is
\[
 \sum_{j=0}^{\lfloor B\rfloor}[x^j]\prod_w(1+x^{c_w}),
\]
interpreted as zero when $B<0$. Thus even the full count is determined
by a histogram of signed occurrence multiplicities, not by a claim of
independence. The pair counts give the simpler sufficient estimate above.

\section{A deterministic counting version of the full proof}

Fix any simple signed interaction graph, with one clause per edge and
no initial units, and fix an order and a depth $T=\lfloor\tau n\rfloor$,
where $0<\tau<1$ is constant. Let $d\ge1$ be a reference degree scale,
$h_i$ the initial literal-degree gap, and $J_{iw}$ the product of the
literal signs on an edge, zero off edges. Define the ordered interaction
row $r_i(w)=J_{iw}$ for $w>i$, and zero otherwise, for $i\le T$.
Its raw pair counts are
\[
 M_{ij}=\sum_{w>\max(i,j)}J_{iw}J_{jw}
       =\#\{\text{agreeing interaction signs}\}
        -\#\{\text{opposite interaction signs}\}.
\]
Here $M_{ii}$ is the number of later neighbors of $i$. Unlike $C$ in
Section 2, these rows have not been oriented by preferred values.

For a prefix $u\in\{-1,+1\}^T$, define its total gap
\[
 G(u)=\sum_{i\le T}\left|h_i-\sum_{j<i}J_{ij}u_j\right|.
\]
This is defined on the full assignment cube without survival tests or
optimization. Fix constants $K>0$, $H\ge0$, $g>0$, and $\eta>0$.
The three conditions are
\begin{align}
 \sum_{i,j\le T}M_{ij}z_iz_j
       &\le Kd\sum_{i\le T}z_i^2
       &&\text{for all }z\in\mathbb R^T,             \label{eq:rawcounts}\\
 \sum_{i\le T}h_i^2&\le Hnd,                         \label{eq:henergy}\\
 \#\{u\in\{-1,+1\}^T:G(u)<gn\sqrt d\}
       &\le 2^T e^{-\eta n}.                        \label{eq:gapcount}
\end{align}
The first is an aggregate condition on pair counts, not a bound on
each pair separately. It can be checked by linear algebra on $M$; the
stronger condition $\sum_j|M_{ij}|\le Kd$ for every $i$ also suffices.
The last condition counts exceptional gap profiles. It is not asserted
to have a polynomial-time membership test. Replacing it by simpler input
conditions is a remaining structural question.

\begin{theorem}[Full deterministic counting criterion]\label{thm:full}
Every fixed formula satisfying~\eqref{eq:rawcounts}--\eqref{eq:gapcount}
is solved exactly by fixed-order branching through depth $T$, followed
by enumeration of the summed learned inequality, in polynomial space and
time
\[
 \Ostar(2^T+2^n e^{-cn}),\qquad
 c=\min\left\{\eta,\frac{g^2}{32(H+K)}\right\}>0.
\]
All constants are independent of $d$. The statement has no probability
qualification over input formulas.
\end{theorem}
\begin{proof}
Use a uniform full assignment $u$ only to count generated candidates.
Let $I$ consist of its nonpreferred decisions among the first $T$,
$v_i=u_i[i\in I]$, and $\Delta_i=h_i-\sum_{j<i}J_{ij}u_j$.
As in the early-enumeration note, set
\[
 S=\sum_{i\in I}(|\Delta_i|/2+1),\qquad
 Q=\frac12\sum_{w=1}^n u_w\sum_{\substack{i<w\\i\le T}}J_{iw}v_i.
\]
Exactly those assignments with $S+Q\le0$ pass their prefix's summed
inequality. Consequently, their total number is
$W(F)=2^n\Prb_u[S+Q\le0]$ for this fixed formula.

The coefficient of $u_w$ in the inner sum for $Q$ is determined before
$u_w$ is exposed. Its total squared coefficient is at most
$Kd\sum_iv_i^2\le KdT$ by~\eqref{eq:rawcounts}. The conditional
cosh bound, compensated by these predictable squared coefficients,
therefore gives
\begin{equation}
 \Prb[Q\le-b]\le e^{-2b^2/(KdT)}.                  \label{eq:fullQ}
\end{equation}

Writing $R$ for the matrix of rows $r_i$, condition~\eqref{eq:rawcounts}
is $\|R\|^2\le Kd$. Its leading $T\times T$ submatrix is the transpose
of the map from the prefix to its gap changes. Hence, for every prefix,
\[
 \sum_{i\le T}\Delta_i^2
 \le2\sum_{i\le T}h_i^2+2KdT\le2(H+K)nd.
\]
Conditional on earlier decisions, $X_i=[i\in I]|\Delta_i|/2$ has mean
$|\Delta_i|/4$ and second moment $\Delta_i^2/8$. Using
$e^{-\lambda X_i}\le1-\lambda X_i+\lambda^2X_i^2/2$ and the predictable
second-moment sum at most $(H+K)nd/4$ gives
\begin{equation}
 \Prb\left[\sum_{i\le T}X_i<G(u)/4-b\right]
       \le e^{-2b^2/((H+K)nd)}.                     \label{eq:fullS}
\end{equation}
These conditional arguments require independent assignment coordinates,
which the counting experiment supplies; no input signs are random.

Take $b=gn\sqrt d/8$ in~\eqref{eq:fullQ}--\eqref{eq:fullS}.
Outside the gap exception in~\eqref{eq:gapcount} and the shortfall
exception in~\eqref{eq:fullS}, $S\ge\sum_iX_i\ge gn\sqrt d/8$.
If also $Q>-gn\sqrt d/8$, then $S+Q>0$. A union bound therefore gives
\[
 \frac{W(F)}{2^n}\le
 e^{-\eta n}+e^{-g^2n/(32(H+K))}+e^{-g^2n/(32K)}
 \le3e^{-cn},
\]
where $T\le n$ was used in the $Q$ bound. The cost is
$\Ostar(2^T+W(F))$, proving the result.
\end{proof}

Thus there is a precise counting replacement for the two roles of
independence. Pair counts control collective agreement even after
adaptive selection and orientation of rows. The gap-profile count
controls how frequently the order fails to produce useful imbalance.
The latter condition is a sufficient structural counting hypothesis,
not a restatement that few assignments survive the algorithm.

\section{A transparent bipartite special case}

Let the formula have a simple bipartite interaction graph with parts
$L$ and $R$, of sizes $p$ and $q$. There is one clause per edge and no
initial unit clauses. Every vertex of $L$ has degree $d$; degrees in
$R$ need not be equal. Let $h_i$ be the positive literal degree minus
the negative literal degree of $i\in L$, and let $t_i\in\{-1,+1\}$
be its preferred value, using $+1$ to break a tie.

For edge $iw$, let $J_{iw}$ be the product of its two literal signs.
The learned inequality in the nonpreferred branch of $i$ counts the
bad value
\begin{equation}
 b_i(w)=-t_iJ_{iw},\qquad
 L_i(y)\le\frac{d-|h_i|}{2}-1.                       \label{eq:learn}
\end{equation}
Set $b_i(w)=0$ off the support. All these rows and their pair counts
$C_{ij}$ are defined directly from the input. They do not depend on
which branches are later taken, because variables in $L$ have no edges
to one another.

Fix constants $a>0$ and $K\ge1$. Impose the following two conditions:
\begin{align}
 |h_i|&\ge a\sqrt d &&\text{for every }i\in L,       \label{eq:gap}\\
 d|I|+2\sum_{\substack{i<j\\i,j\in I}}C_{ij}
      &\le Kd|I| &&\text{for every }I\subseteq L.   \label{eq:counts}
\end{align}
Condition~\eqref{eq:counts} bounds the total excess of agreements over
disagreements in every selected family. It permits nonzero correlations
and arbitrary higher-order dependencies.

One stronger, polynomial-time checkable condition is
\begin{equation}
 \sum_{j\ne i}\max(C_{ij},0)\le(K-1)d
       \quad\text{for every }i\in L.               \label{eq:rowsum}
\end{equation}
Indeed, the off-diagonal sum in~\eqref{eq:counts} is at most the sum of
these positive row sums over $i\in I$. Another sufficient certificate is
$C\preceq Kd\,\mathrm{Id}$, which can be checked by linear algebra.
That certificate is an equivalent spectral formulation of a stronger
quadratic bound; it is not a separate independence assumption. The
proof needs only the subset inequality~\eqref{eq:counts}.

\begin{theorem}[Counting conditions suffice]\label{thm:det}
Every input satisfying~\eqref{eq:gap}--\eqref{eq:counts} can be solved
exactly in polynomial space and time
\begin{equation}
 \Ostar\left(2^p+2^q(1+e^{-a^2/(2K)})^p\right).     \label{eq:time}
\end{equation}
If $p/(p+q)$ is bounded away from zero and one, this is
$\Ostar((2-\delta)^{p+q})$ for a fixed $\delta>0$ independent of $d$.
The stronger condition~\eqref{eq:rowsum}, together with
\eqref{eq:gap}, defines a class whose membership can be verified in
polynomial time when the bipartition is supplied.
\end{theorem}
\begin{proof}
Branch on every variable of $L$, learning~\eqref{eq:learn} when its
nonpreferred value is used. At each resulting prefix freeze the learned
family and enumerate solutions of its summed inequality. The learned
constraints preserve the first optimum in preferred-child-first order:
violating a learned constraint permits a no-worse single flip to the
preferred value, producing an earlier optimum, a contradiction.
Their sum therefore also preserves an optimum.

There is exactly one prefix for each set $I$ of nonpreferred decisions.
For $k=|I|>0$, its shortfall and variance parameter satisfy
\[
 s_I=\sum_{i\in I}(|h_i|/2+1)\ge ka\sqrt d/2,
 \qquad V_I\le Kdk.
\]
Lemma~\ref{lem:count} bounds its candidate count by
$N(I)\le2^q e^{-a^2k/(2K)}$. The same formula holds for $I=\varnothing$,
whose count is $2^q$. Summing over all prefixes gives
\[
 \sum_{I\subseteq L}N(I)
 \le2^q\sum_{k=0}^p\binom pk e^{-a^2k/(2K)}
 =2^q(1+e^{-a^2/(2K)})^p.
\]
After orienting cheaper values, a depth-first generator of the
nonnegative budget inequality takes $\operatorname{poly}(n)(N(I)+1)$
time: every retained partial assignment extends by choosing cheaper
values on the remaining variables. This proves~\eqref{eq:time}; the
traversal uses polynomial space. If $p=\tau n$, the two exponential
bases are $2^\tau$ and
$2((1+e^{-a^2/(2K)})/2)^\tau$, both uniformly below two when $\tau$
stays in a compact subinterval of $(0,1)$.
\end{proof}

This special case is a theorem about the constraint-guided enumeration
mechanism, not a new tractability result for bipartite Max 2-SAT.
Every bipartite formula already has an $\Ostar(2^{\min(p,q)})$ algorithm:
enumerate the smaller side and optimize the variables on the other
side independently. The value of the special case is its direct
translation from static neighborhood counts to a bound on every
selected learned family. Theorem~\ref{thm:full} is not restricted to
bipartite graphs or graphs with a large independent set.

\section{An explicit dense example}

Take $p=q=d=4^k$ for any integer $k\ge1$, with the complete bipartite
interaction graph. Let $H$ be a $q\times q$ sign matrix satisfying
$HH^{\mathsf T}=q\,\mathrm{Id}$. Such matrices follow from the recursion
\[
 H_1=(1),\qquad
 H_{2r}=\begin{pmatrix}H_r&H_r\\H_r&-H_r\end{pmatrix};
\]
the stated identity follows by block multiplication.
Choose $u\in\{-1,+1\}^q$ with $\sum_wu_w=\sqrt q$.
For each edge $iw$, set the literal sign at $i$ to $u_w$ and that
at $w$ to $u_wH_{iw}$. This specifies all input clauses explicitly.

Every left gap is $\sqrt q$, the interaction sign matrix is $H$, and
the bad-value matrix is $-H$. Consequently $C_{ij}=0$ for $i\ne j$:
each pair agrees on exactly half its neighbors and disagrees on the
other half. Conditions~\eqref{eq:gap}--\eqref{eq:counts} hold with
$a=K=1$. Since $n=2q$, Theorem~\ref{thm:det} gives
\[
 \Ostar\left(\left[\sqrt{2(1+e^{-1/2})}\right]^n\right)
 =\Ostar(1.792502^n)\subseteq\Ostar(1.793^n).
\]
The clause density is $q/2=n/4$. There is no random-input qualification.
This is an explicit witness that the class is nonempty at growing
density, not a claim that this structured example requires exponential
time or that the constants are optimal.

\section{Why cancellation counts alone do not suffice}

The role of the gap hypothesis is real for this particular algorithm.
Take a matrix with $p$ orthogonal sign rows of length $q$, where
$q\ge16p$, and make it the interaction matrix of a complete bipartite
formula. Choose the own-literal signs at each left variable to be
balanced, and choose the other endpoint signs to realize the desired
interaction products. Then every $h_i=0$, while every distinct pair
has $C_{ij}=0$ exactly. For example, take the first $p=q/16$ rows of
the recursively defined matrix above, for powers of two $q\ge32$.
These instances have clause density proportional to their number of
variables.

\begin{proposition}[Failure of the summed generator with zero gaps]
For these formulas, the algorithm that branches on $L$ and then
enumerates the summed inequality generates at least
\[
 \frac3{32}(2^p-1)2^q
\]
complete assignments. Thus its work is $\Omega(2^{p+q})$ despite
perfectly balanced signed pair-agreement counts.
\end{proposition}
\begin{proof}
For any nonempty set $I$ of $k$ nonpreferred decisions, the total
shortfall is $s_I=k$ and $V_I=qk$. Under a uniform right assignment,
$Q_I$ is a symmetric sum of independent signs, of variance
$v=qk/4$ and fourth moment at most $3v^2$. For completeness, the
elementary second-moment inequality for $X\ge0$ is
\[
 \Prb[X\ge\theta\E X]\ge
 (1-\theta)^2\frac{(\E X)^2}{\E X^2},\qquad 0\le\theta<1.
\]
It follows by bounding the contribution from $X<\theta\E X$ and
applying Cauchy--Schwarz to the remaining contribution. Apply it to
$X=Q_I^2$ and $\theta=k^2/v=4k/q\le1/4$. Then
\[
 \Prb[|Q_I|\ge k]\ge\frac{(1-4k/q)^2}{3}\ge\frac3{16}.
\]
Symmetry and $k>0$ give
$\Prb[Q_I\le-k]\ge3/32$. These are exactly the completions admitted
by the summed inequality. Sum their counts over nonempty $I$.
\end{proof}

This does not show that the formulas are hard for other branching
orders, domination, individual-constraint propagation, or aggregation.
It shows that initial literal balance and pairwise neighborhood balance
do not, by themselves, give the earlier switching guarantee for this
generator. One must also obtain sufficiently restrictive constraints.

\clearpage
\section{What remains for the general branching approach}

For a fixed residual constraint family, equations~\eqref{eq:C}--\eqref{eq:V}
and Lemma~\ref{lem:count} already apply to arbitrary inputs. The
remaining global task is to guarantee a large shortfall relative to
$V_I$ on sufficiently many prefixes, accounting for the cost of all
other prefixes. In the bipartite theorem, static gaps and the binomial
count of nonpreferred decisions settle this exactly. In a general
formula, branching changes gaps, supports, and which rows are selected.
Initial pair counts cannot simply be reused after these restrictions.

This gives a concrete form of the proposed dichotomy. Either learning
creates a large total shortfall while the selected neighborhoods have
controlled net agreement, and enumeration saves an exponential factor;
or one must address insufficient gaps or a selected family with large
positive agreement excess. The latter is a precisely countable alignment
of bad values, but it is not yet a proof of a small separator or an
aggregation rule with a global running-time bound.

The counts also give exact transitions for a fixed summed inequality.
Write its current shortfall as $s$ and let
$B_w=\sum_i b_i(w)$, so $V=\sum_wB_w^2$. Assigning $y_w=\varepsilon$
changes these parameters to
\[
 s'=s+\frac{\varepsilon B_w}{2},\qquad V'=V-B_w^2.
\]
Thus the effect of choosing a variable can be read from occurrence
counts directly. A large $|B_w|$ gives a large separation between the
two children's shortfalls and removes a large contribution to $V$.
This does not itself shrink the number of completions of the inherited
inequality: their counts partition between the children. Further savings
must come from additional constraints, domination, or another valid
optimization argument.

\paragraph{Verification.}
The companion program \texttt{explore\_counting\_conditions.py} verifies
the agreement identities, the upper bound and preservation of an optimum
on small explicit instances, and the zero-gap lower bound using exact
completion counts. All 548 prefix checks passed, including 65,792
full-assignment identity checks on the small positive-gap examples.
Another 2,560 complete-assignment checks on signed complete graphs
verified the ordered-count identities and quadratic bounds used in
Theorem~\ref{thm:full}. These checks support the implementation; the
theorems follow from the proofs above.

\paragraph{Related working note.}
The algorithm and output-sensitive budget generator are those of
\emph{Max 2-SAT by Early Constraint-Guided Enumeration}, 8 October 2026,
local source \texttt{max2sat\_early\_enumeration.tex}. The present note
replaces a probabilistic assumption by sufficient deterministic
conditions for a narrower input class. It does not assert that the
conditions are necessary or characterize typical random signings.
\end{document}
